% from ega4-18-fr.tex
\oldpage[IV-4]{211}

  \item[f$'$)] L'homomorphisme $\operatorname{gr}_\bullet(R)\otimes_A(A/\mathfrak J)\to\operatorname{gr}_\bullet(N)\otimes_A(A/\mathfrak J)$ obtenu en restreignant $j'$ aux polynômes de degré $0$, est injectif.
  \item[g)] La filtration $\mathfrak J$-préadique de $\operatorname{gr}_\bullet(R)$ est induite par celle de $\operatorname{gr}_\bullet(N)$.
  \item[h)] Les filtrations $(R_k')$ et $(R_k'')$ sur $R$ sont identiques.
\end{enumerate}
On a alors les implications suivantes:
\[
\xymatrix@C=2.4em@R=1.8em{
  a) \ar@{=>}[r] \ar@{=>}[dr] & e) \ar@{=>}[r] \ar@{=>}[dr] & b) \ar@{<=>}[r] & c) \ar@{<=>}[r] & h) \\
  & g) \ar@{=>}[r] & d) \ar@{<=>}[r] & d') \ar@{<=>}[r] & f) \ar@{<=>}[r] & f')
}
\]
En outre, lorsque tout module quotient d'un sous-module d'un $\operatorname{gr}_k(M)$ est séparé pour la topologie $\mathfrak J$-préadique (ce qui sera le cas lorsque $A$ est noethérien, que les $f_i$ appartiennent au radical de $A$ et que les $\operatorname{gr}_k(M)$ sont des $A$-modules de type fini), alors les conditions $a)$ à $h)$ sont toutes équivalentes.
\end{proposition}

Notons que d'après (19.5.5), de l'hypothèse que $\mathbf f$ est $\operatorname{gr}_\bullet(N)$-régulière, il résulte que $\psi_N$ est \emph{bijectif}; comme $\operatorname{gr}'(j)$ est injectif, l'équivalence de $d)$ et $f)$ résulte du diagramme (19.6.2.1), ainsi que celle de $d')$ et $f')$; par ailleurs $f)$ et $f')$ sont trivialement équivalentes, ce qui montre l'équivalence de $d)$, $d')$, $f)$ et $f')$.

Comme on sait déjà que $\psi_M$ est surjectif et $\psi_N$ bijectif, l'équivalence de $b)$ et $c)$ découle aussi du diagramme (19.6.2.1) par un cas particulier du lemme des cinq.

La relation $\psi_R'=\gamma\circ\psi_R$ (19.6.2) et le fait que $\psi_R$ est surjectif montrent que la condition $c)$ équivaut à dire que $\gamma$ est surjectif, ce qui équivaut encore à la condition $h)$; on a donc prouvé l'équivalence de $b)$, $c)$ et $h)$.

Il est trivial que $e)$ équivaut à la conjonction de $c)$ et $d)$, donc aussi de $b)$ et $d)$. Notons maintenant que $a)$ implique $b)$ et que $b)$ implique $a)$ sous l'hypothèse de séparation (19.5.5). D'autre part, les implications $a)\Rightarrow g)\Rightarrow f')$, et l'implication $f')\Rightarrow a)$ sous l'hypothèse de séparation, résultent de (19.5.3). On a donc prouvé que $a)$ entraîne $e)$ (équivalente à la conjonction de $b)$ et $f')$), et aussi que toutes les conditions sont équivalentes sous l'hypothèse de séparation. C.Q.F.D.

\begin{corollary}[19.6.4]
\label{IV.19.6.4}
Soient $\mathfrak R$ un idéal de $A$, $\mathfrak L=\mathfrak J+\mathfrak R$. Supposons que la filtration $(N_k)$ (donc aussi $(M_k)$) soit la filtration $\mathfrak R$-préadique, de sorte que les filtrations $(M_k')$ et $(N_k')$ sont les filtrations $\mathfrak L$-préadiques. Considérons $\operatorname{gr}_\bullet(R)$ (resp. $\operatorname{gr}_\bullet'(R)$) comme un module gradué sur $\operatorname{gr}_{\mathfrak R}^*(A)$ (resp. $\operatorname{gr}_{\mathfrak L}^*(A)$).
\begin{enumerate}
  \item[(i)] Soit $S$ une partie de $\operatorname{gr}_\bullet(R)$ qui engendre $\operatorname{gr}_\bullet(R)$ comme $\operatorname{gr}_{\mathfrak R}^*(A)$-module. Alors la condition $c)$ de (19.6.3) équivaut à la condition suivante:
  \item[c$'$)] L'image de $S$ dans $\operatorname{gr}_\bullet'(R)$ par $\psi_R'$ engendre $\operatorname{gr}_\bullet'(R)$ comme $\operatorname{gr}_{\mathfrak L}^*(A)$-module.
  \item[(ii)] Supposons vérifiée la condition $e)$ de (19.6.3), et de plus que les quotients des $\operatorname{gr}_k(R)$ soient séparés pour la topologie $\mathfrak J$-préadique (ce qui sera le cas lorsque $A$ est noethérien, les $f_i$ dans le radical de $A$ et les $\operatorname{gr}_k(N)$ des $A$-modules de type fini). Alors, pour qu'une partie $S$ de $\operatorname{gr}_\bullet(R)$ formée d'éléments homogènes engendre $\operatorname{gr}_\bullet(R)$ comme $\operatorname{gr}_{\mathfrak R}^*(A)$-module, il faut et il suffit que son image par $\psi_R'$ engendre $\operatorname{gr}_\bullet'(R)$ comme $\operatorname{gr}_{\mathfrak L}^*(A)$-module.
\end{enumerate}
\end{corollary}

